\SourceLeaf{401}%
% A544-P0401-D01
\[
a^7+b^7+c^7=0
\]
nothwendig eine der Grössen $a,b,c$ durch $7$ theilbar sein müsse. Es folgt nemlich nur
\[
\begin{aligned}
(\alpha+\beta+\gamma)^7={}&5\alpha\beta\gamma\{3(\alpha^4+\beta^4+\gamma^4)\\
&\qquad+10(\alpha\alpha\beta\beta+\alpha\alpha\gamma\gamma+\beta\beta\gamma\gamma)\}
\end{aligned}
\]
welches bestehen kann, ohne dass $\alpha,\beta,\gamma$ durch $1-\varepsilon$ theilbar wäre.

% A544-P0401-D02
Hoffentlich wird sich indessen dies in Zukunft aus der Natur der Determinanten und der Einheitszahlen ableiten lassen.
\BookRule{.3\textwidth}

% A544-P0401-D03
\pdfbookmark[1]{II. Bestimmung der nächsten ganzen Zahl}{complexe-II}
\begin{center}[II.]\par\medskip
BESTIMMUNG DER NACHSTEN GANZEN ZAHL.
\end{center}
Es sei
\[
\begin{aligned}
\varepsilon^3&=1,\quad m=a+b\varepsilon+c\varepsilon\varepsilon\\
2a-b-c&=A+\alpha\\
2b-c-a&=B+\beta\\
2c-a-b&=C+\gamma
\end{aligned}
\]
wo $A,B,C$ ganze Zahlen; $\alpha,\beta,\gamma$ positive echte Brüche sind. Man hat dann
\[
A+B+C+\alpha+\beta+\gamma=0
\]
also drei Fälle zu unterscheiden:

% A544-P0401-D04
I. $\alpha+\beta+\gamma=0$, folglich $\alpha=0$, $\beta=0$, $\gamma=0$

1, $A\equiv B\equiv C\ (\text{mod. }3)$. Hier ist $m$ selbst eine ganze Zahl.

2, $A-B\equiv B-C\equiv C-A\equiv\pm1\ (\text{mod. }3)$. Hier ist $m\pm\frac{\varepsilon-\varepsilon\varepsilon}{3}\pdot\varepsilon^n$ eine ganze Zahl.

% A544-P0401-D05
II. $\alpha+\beta+\gamma=1$.

Hier ist
\[
\begin{aligned}
&A+B\varepsilon+C\varepsilon\varepsilon+1\\
&A+B\varepsilon+C\varepsilon\varepsilon+\varepsilon\\
&A+B\varepsilon+C\varepsilon\varepsilon+\varepsilon\varepsilon
\end{aligned}
\]
\SourceLeaf{402}%
% A544-P0402-D01
jedes durch $1-\varepsilon$ theilbar, und eine dieser Zahlen durch $3$. Der Quotient oder
\[
m+\frac{\varepsilon^n-\alpha-\beta\varepsilon-\gamma\varepsilon\varepsilon}{3}
\]
die gesuchte ganze Zahl.

% A544-P0402-D02
III. $\alpha+\beta+\gamma=2$

Hier sind
\[
\begin{aligned}
&A+B\varepsilon+C\varepsilon\varepsilon+\varepsilon+\varepsilon\varepsilon\\
&A+B\varepsilon+C\varepsilon\varepsilon+\varepsilon\varepsilon+1\\
&A+B\varepsilon+C\varepsilon\varepsilon+1+\varepsilon
\end{aligned}
\]
durch $1-\varepsilon$ und eine dieser Zahlen durch $3$ theilbar. Der Quotient, oder
\[
m+\frac{\varepsilon^n(\varepsilon+\varepsilon\varepsilon)-\alpha-\beta\varepsilon-\gamma\varepsilon\varepsilon}{3}
\]
ist die gesuchte ganze Zahl.

% A544-P0402-D03
In allen drei Fällen hat der Rest die Form
\[
x+y\varepsilon+z\varepsilon\varepsilon
\]
so dass $x,y,z$ ohne Rücksicht auf das Zeichen kleiner als $\frac13$ und $x+y+z=0$ wird. Dadurch wird aber nothwendig
\[
xx+yy+zz=2xx-2yz=2yy-2xz=2zz-2xy<\tfrac29
\]
weil von den drei Grössen $x,y,z$ nothwendig zwei einerlei Zeichen haben. Folglich ist der Determinant des Restes
\[
=\tfrac32(xx+yy+zz)<\tfrac13\qquad\text{Q. E. D.}
\]

% A544-P0402-D04
Die Bestimmung der nächsten ganzen Zahl geschieht \emph{bequemer} auf folgende Art. Es sei vorgegeben $a+b\varepsilon+c\varepsilon\varepsilon=m$, man setze
\[
\begin{aligned}
b-a&=C+\gamma\\
c-b&=A+\alpha\\
a-c&=B+\beta
\end{aligned}
\]
wo $A,B,C$ die nächst kleinern ganzen Zahlen; $\alpha,\beta,\gamma$ positive Brüche sind.
Hier sind drei Fälle zu unterscheiden:

\SourceLeaf{403}%
% A544-P0403-D01
I. $\alpha+\beta+\gamma=0$, so ist $m$ selbst ganze Zahl

II. $\alpha+\beta+\gamma=1$, so ist die nächste ganze Zahl
\[
\begin{array}{ll}
B+(B+C)\varepsilon&\text{wenn }\alpha\text{ der grösste Bruch ist.}\\
\qquad C\varepsilon+(A+C)\varepsilon\varepsilon&\qquad\beta\\
A+B\quad\pdot\quad+A\varepsilon\varepsilon&\qquad\gamma
\end{array}
\]

% A544-P0403-D02
II. $\alpha+\beta+\gamma=2$, so ist die nächste ganze Zahl
\[
\begin{array}{ll}
B+1+(B+C+2)\varepsilon&\text{wenn }\alpha\text{ der kleinste Bruch ist.}\\
\qquad(C+1)\varepsilon+(A+C+2)\varepsilon\varepsilon&\qquad\beta\\
A+B+2\quad\pdot\quad+(A+1)\varepsilon\varepsilon&\qquad\gamma
\end{array}
\]
In II, $1$ ist der Rest $\beta+(\beta+\gamma)\varepsilon$, dessen Determinant
\[
\begin{aligned}
&=\beta\beta+\beta\gamma+\gamma\gamma\\
&=\tfrac13-\tfrac16[(\alpha-\beta)(1+3\beta)+(\alpha-\gamma)(1+3\gamma)]
\end{aligned}
\]

% A544-P0403-D03
\emph{Noch einfacher so:}

Man ordne die Brüche $a-[a],b-[b],c-[c]$ nach ihrer Grösse: so heissen sie der Reihe nach $p,q,r$. Sind alle drei gleich gross, so ist $m$ eine ganze Zahl. Sind sie aber ungleich, so sei $t$ ein beliebiger Bruch zwischen
\[
\begin{array}{llll}
p\text{ und }q,&\text{jenachdem}&q-p&\text{am grössten ist}\\
q\text{ und }r&&r-q&\\
r\text{ und }1+p&&1+p-r&
\end{array}
\]
Sodann ist
\[
[a-t]+[b-t]\varepsilon+[c-t]\varepsilon\varepsilon
\]
die nächste ganze Zahl.
\BookRule{.3\textwidth}

\SourceLeaf{404}%
% A544-P0404-D01
\pdfbookmark[1]{III.}{complexe-III}
\begin{center}[III.]\end{center}
Es sei $\varepsilon^5=1$
\[
\begin{aligned}
a+b\varepsilon+c\varepsilon\varepsilon+d\varepsilon^2+e\varepsilon^4&=q'\\
a+b\varepsilon^{-1}+c\varepsilon^{-2}+d\varepsilon^{-3}+e\varepsilon^{-4}&=q''''
\end{aligned}
\]
\[
\begin{aligned}
(a-b)^2+(b-c)^2+(c-d)^2+(d-e)^2+(e-a)^2&=2p'\\
(a-c)^2+(b-d)^2+(c-e)^2+(d-a)^2+(e-b)^2&=2p''
\end{aligned}
\]
\[
\begin{aligned}
q'q''''&=-p'\varepsilon-p''\varepsilon\varepsilon-p''\varepsilon^3-p'\varepsilon^4=P'\\
q''q'''&=-p'\varepsilon\varepsilon-p''\varepsilon^4-p''\varepsilon-p'\varepsilon^3=P''
\end{aligned}
\]

% A544-P0404-D02
\[
\text{Determinant}=P'P''=-p'p'+3p'p''-p''p''
\]
\[
\begin{aligned}
\text{Mensura}&=2p'+2p''=2P'+2P''\\
&=5(aa+bb+cc+dd+ee)-(a+b+c+d+e)^2
\end{aligned}
\]
Multiplicando per $1-\varepsilon$ fit mensura nova $=8p'$
\[
\text{Höchste Mensur}=2\left(\frac{\sin72^\circ}{\sin36^\circ}+\frac{\sin36^\circ}{\sin72^\circ}\right)\surd D=4{,}472\surd D
\]

% A544-P0404-D03
\emph{Modulus} $=1-\varepsilon$
\[
\begin{aligned}
1-\varepsilon&=x\\
\varepsilon&=1-x\\
\varepsilon\varepsilon&=1-2x+xx\\
\varepsilon^3&=1-3x+3xx-x^3\\
\varepsilon^4&=1-4x+6xx-4x^3+x^4\\
&=-4+6x-4xx+x^3
\end{aligned}
\]
Also
\[
\begin{aligned}
\frac{1-\varepsilon^n}{1-\varepsilon}&\equiv n\text{ mod. }(1-\varepsilon)\\
\varepsilon^n&\equiv1-nx\text{ mod. }(1-\varepsilon)^2\\
\left(\frac{\varepsilon+\varepsilon^4}{\varepsilon\varepsilon+\varepsilon^3}\right)^n&\equiv1+nxx\text{ mod. }(1-\varepsilon)^3
\end{aligned}
\]
Also eine Zahl, welche $\equiv1\text{ mod. }(1-\varepsilon)^3$ kann \emph{nur} dann eine Einzahl sein, wenn sie zugleich $\equiv1\ (\text{mod. }5)$.
\BookRule{.3\textwidth}

\SourceLeaf{405}%
% A544-P0405-D01
\pdfbookmark[1]{IV. Einiges über die Mensur der Zahlen}{complexe-IV}
\begin{center}[IV.]\par\medskip
EINIGES UBER DIE MENSUR DER ZAHLEN.
\end{center}
Es sei $\varepsilon^n=1$, $n$ Primzahl
\[
\begin{aligned}
m&=a+a'\varepsilon+a''\varepsilon\varepsilon+a'''\varepsilon^3+\ldots+a^{(n-1)}\varepsilon^{n-1}=f\varepsilon\\
D&=f\varepsilon\pdot f\varepsilon\varepsilon\pdot f\varepsilon^3\ldots f\varepsilon^{n-1}\\
f\varepsilon\pdot f\varepsilon^{n-1}&=-b'(\varepsilon+\varepsilon^{n-1})-b''(\varepsilon^2+\varepsilon^{-2})\\
&\qquad-b'''(\varepsilon^3+\varepsilon^{-3})\ldots
\end{aligned}
\]
so ist
\[
\begin{aligned}
2b'&=(a-a')^2+(a'-a'')^2+(a''-a''')^2+\text{etc.}\\
2b''&=(a-a'')^2+(a'-a''')^2+(a''-a'''')^2+\text{etc.}\\
&\text{etc.}
\end{aligned}
\]

% A544-P0405-D02
hier sind also $b',b'',b'''\dotstwo$ lauter positive Grössen; sie heissen \emph{Partialmensuren} von $m$, so wie ihre Summe
\[
\begin{aligned}
b'+b''+b'''+\text{etc.}={}&n(aa+a'a'+a''a''+\dotstwo)\\
&\qquad-(a+a'+a''+\text{etc.})^2
\end{aligned}
\]
die \emph{Generalmensur}. Setzt man
\[
f\varepsilon\pdot f\varepsilon^{n-1}=c',\quad f\varepsilon\varepsilon\pdot f\varepsilon^{n-2}=c''\quad\text{etc.}
\]
so ist
\[
\begin{aligned}
&c'+c''+c'''+\text{etc.}+c^{\frac12(n-1)}\\
&\qquad=b'+b''+b'''+\text{etc.}+b^{\frac12(n-1)}
\end{aligned}
\]
\[
\begin{aligned}
&c'(\varepsilon+\varepsilon^{n-1})+c''(\varepsilon\varepsilon+\varepsilon^{n-2})+c'''(\varepsilon^3+\varepsilon^{n-3})+\text{etc.}\\
&\qquad=2(b'+b''+b'''+\text{etc.}+b^{\frac12(n-1)})-nb'
\end{aligned}
\]
\[
\begin{aligned}
&c'(2-\varepsilon-\varepsilon^{n-1})+c''(2-\varepsilon\varepsilon-\varepsilon^{n-2})\\
&\qquad+c'''(2-\varepsilon^3-\varepsilon^{n-3})+\text{etc.}=nb'
\end{aligned}
\]
\[
b'>\frac{n-1}{2n}(nD)^{\frac2{n-1}},\quad b'+b''+b'''+\text{etc.}>\frac{n-1}{2}\pdot D^{\frac2{n-1}}
\]

% A544-P0405-D03
Ist allgemein
\[
f\varepsilon\pdot f\varepsilon^{n-1}=A+A'\varepsilon+A''\varepsilon\varepsilon+A'''\varepsilon^3+\dotstwo
\]
so ist die Generalmensur $\triangle=-A-A'-A''-\text{etc.}+nA$

Mensur von $(1+\varepsilon)f\varepsilon\ldots$
\[
\triangle'=4\triangle-2n(A-A')=4\triangle-2nb'
\]
Ist $a+a'+a''+\dotstwo=0$, so ist $\triangle=n(aa+a'a'+a''a''+\text{etc.})$

und ist $A+A'+A''+\ldots=0$, so ist $\triangle=nA$, $\triangle'=n(2A+2A')$

% A544-P0405-D04
Ist also einer der Coëfficienten $A',A''$ etc. negativ und absolut grösser als $\frac12A$, so lässt sich die Mensur salvo determinante herabbringen.
\BookRule{.3\textwidth}

